\section{Results}
\label{sec:results}

\subsection{Release Integrity}

The experiment completed and the release passed a fresh-root restore. All 11
Stage-09 artifacts matched their declared SHA-256 hashes and byte counts. The clean
set contains 2,040 unique units and the robustness set 1,440; both have zero
duplicate IDs, divergent route hashes, non-finite audited values, or execution
failures. Generation-window energy coverage is 100\% in both sets
(Table~\ref{tab:trace-integrity}). Thus the negative outcome is not attributable to
missing runs or telemetry filtering.

\begin{table*}[t]
\caption{Released trace integrity and outcome audit. Both sets had zero duplicate unit IDs, zero non-finite audited values, zero divergent route hashes, and zero execution failures.}
\label{tab:trace-integrity}
\centering
\small
\begin{tabular}{@{}lrrrrrr@{}}
\toprule
Set & Rows & Routes & Strict successes & Parse & Support & Energy coverage \\
\midrule
clean & 2040 & 17 & 229 (11.2\%) & 58.4\% & 18.0\% & 100.0\% \\
robustness & 1440 & 3 & 60 (4.2\%) & 49.8\% & 6.7\% & 100.0\% \\
\bottomrule
\end{tabular}
\end{table*}


Across all clean routes, 58.38\% of outputs parsed, 19.41\% passed the answer check,
18.04\% passed the support guard, and 11.23\% met every strict condition. These
aggregate rates are diagnostic rather than model-weighted benchmark scores because
the route matrix deliberately repeats every question under many mechanisms.

\subsection{Confirmatory Router Outcome}

\begin{table}[t]
\caption{Confirmatory policy result on the sealed test split ($n=120$). Intervals are 95\% scenario-cluster bootstrap intervals.}
\label{tab:confirmatory}
\centering
\small
\begin{tabular}{@{}lr@{}}
\toprule
Quantity & Result \\
\midrule
Selected route & A03 for all queries \\
Strict success & 0.0\% [0.0, 0.0] \\
Generation-window GPU energy & 143.16 J/query \\
Energy interval & [138.14, 148.07] J \\
Policy latency component sum & 3.76 s/query \\
Abstention / execution coverage & 15.0\% / 100.0\% \\
Policy $-$ baseline success & 0.0 pp \\
Policy $-$ baseline energy & 0.00 J/query \\
\midrule
Formal margin check & pass; policy $=$ baseline \\
Energy reduction & \textbf{fail} \\
Coverage and abstention gates & pass \\
Joint confirmatory hypothesis & \textbf{fail} \\
\bottomrule
\end{tabular}
\end{table}


The frozen policy invoked A03 for all 120 questions. Its strict success was 0\%
(
\ci{} [0, 0]), mean generation-window GPU energy was 143.16 J/query
([138.14, 148.07]), and the summed policy-latency components averaged 3.760 s/query
([3.652, 3.863]). Execution coverage was 100\% and abstention 15\%.

The fixed baseline was also A03. Consequently, both paired differences were exactly
zero: $\Delta_S=0$ and $\Delta_E=0$ with degenerate intervals. The formal relative
quality and operating gates passed, but the upper energy bound was not below zero.
The joint hypothesis therefore failed. The quality result is floor equivalence,
not preservation of useful quality.

Test confidence occupied a single 0--0.1 bin, with mean 0.00937, Brier score
$8.79\times10^{-5}$, and 10-bin ECE 0.00937. Because every selected outcome was
zero, these small errors indicate a correctly pessimistic but non-discriminating
policy; they are not evidence of broadly effective calibration.

\subsection{Route-Level Quality--Energy Structure}

\begin{figure*}[t]
    \centering
    \includegraphics[width=0.93\textwidth]{route_quality_energy.png}
    \caption{Strict success versus generation-window selected-GPU board energy for
    the 17 retained clean routes. Offline reference routes provide the best
    quality points but were not eligible router actions. Cross-board comparisons
    are descriptive; Table~\ref{tab:hardware} gives route assignment.}
    \label{fig:route-scatter}
\end{figure*}

\begin{table}[t]
\caption{Selected clean-test route results. GPU J denotes mean generation-window selected-board energy; the complete matrix and hardware assignment are in Appendix~\ref{app:routes}.}
\label{tab:route-highlights}
\centering
\scriptsize
\begin{tabular}{@{}lrrr@{}}
\toprule
ID & Success (\%) & GPU J & Median s \\
\midrule
M18 & 44.2 & 176.68 & 2.65 \\
A14 & 42.5 & 256.13 & 3.79 \\
M20 & 34.2 & 148.93 & 2.34 \\
M19 & 23.3 & 107.42 & 1.67 \\
M17 & 20.8 & 57.14 & 0.45 \\
A00 & 12.5 & 54.71 & 1.44 \\
A03 & 0.0 & 143.16 & 3.47 \\
\bottomrule
\end{tabular}
\end{table}


Figure~\ref{fig:route-scatter} shows that the best route was the offline
Granite-4.1-3B grounded endpoint M18: 44.17\% success at 176.68 J/query and 2.647-s
median route latency. The grounded Qwen3-4B reference A14 achieved 42.50\% at
256.13 J, while grounded SmolLM3 M20 achieved 34.17\% at 148.93 J. Among direct
references, M19 reached 23.33\%, M17 20.83\%, and M21 13.33\%. The only nonzero
online tiny/small direct endpoint was A00 at 12.50\% and 54.71 J.

The grounded Pareto designation in the release applies only to the five grounded
model endpoints in success--energy--latency space; it identifies the Granite-4.1-3B
and SmolLM3 families. It is not a global frontier over all 17 routes, and the best
grounded points were not deployable policy actions.

\subsection{Generator-Dependent Grounding}

\begin{figure*}[t]
    \centering
    \includegraphics[width=0.94\textwidth]{grounding_effect_forest.png}
    \caption{Matched grounded-minus-direct strict-success effects. Filled markers
    denote intervals that exclude zero; the open SmolLM3 marker spans zero. The
    energy column reports the paired generation-window GPU-energy increase. The
    Granite-4.0-1B energy pair is cross-board and descriptive.}
    \label{fig:grounding-forest}
\end{figure*}

\begin{table*}[t]
\caption{Matched direct-to-grounded transfer on the same 120 sealed-test questions. Success effects are percentage points; energy effects are generation-window selected-GPU board joules per query. ``Same board'' indicates whether both endpoints ran on the same physical T4.}
\label{tab:model-transfer}
\centering
\small
\begin{tabular}{@{}lrrrrr@{}}
\toprule
Generator & Direct (\%) & Grounded (\%) & $\Delta$ success [95\% CI] & $\Delta$ GPU J [95\% CI] & Same board \\
\midrule
Qwen3 0.6B & 12.5 & 0.0 & -12.5 [-18.3, -6.7] & +146.72 [+141.56, +151.62] & yes \\
Granite 4.0 1B & 0.0 & 0.0 & +0.0 [+0.0, +0.0] & +3.88 [+1.73, +6.14] & no \\
Granite 4.1 3B & 20.8 & 44.2 & +23.3 [+10.0, +37.5] & +119.53 [+108.73, +129.74] & yes \\
SmolLM3 3B & 23.3 & 34.2 & +10.8 [-2.5, +24.2] & +41.52 [+35.47, +47.50] & yes \\
Qwen3 4B Instruct & 13.3 & 42.5 & +29.2 [+16.7, +41.7] & +171.71 [+157.74, +185.24] & yes \\
\bottomrule
\end{tabular}
\end{table*}


The matched panel in Fig.~\ref{fig:grounding-forest} and
Table~\ref{tab:model-transfer} rejects a model-independent account of grounding.
Within the equal eight-task mixture, Qwen3-0.6B changed from 12.50\% direct success
to zero grounded success ($-12.50$ \pp{}, interval $[-18.33,-6.67]$). That aggregate
loss came entirely from the no-retrieval stratum: A00 copied all 15 supplied tokens,
whereas M16 copied none. Both routes already scored zero on every evidence-required
stratum. It is therefore inaccurate to claim that RAG reduced this model's success
on questions that needed evidence; the experiment instead shows that an
always-grounded policy destroyed its one functioning direct behavior without
creating evidence-task success.

Granite-4.1-3B improved from 20.83\% to 44.17\% ($+23.33$ \pp{}, [10.00, 37.50]),
and Qwen3-4B improved from 13.33\% to 42.50\% ($+29.17$ \pp{}, [16.67, 41.67]).
SmolLM3 increased from 23.33\% to 34.17\%, but its interval
$[-2.50,24.17]$ includes zero. Positive aggregate effects arose mainly from
knowledge, memory, temporal-update, and authority-conflict tasks. The same grounded
route frequently lost direct-route successes on no-retrieval and deleted/missing
tasks. Hence the estimand is ``always grounded minus always direct'' under the
frozen equal task mixture, not the benefit of retrieving only when retrieval is
needed.

All five grounded endpoints had higher observed generation-window GPU energy. The
same-board increases were 146.72 J for Qwen3-0.6B, 119.53 J for Granite-4.1-3B,
41.52 J for SmolLM3, and 171.71 J for Qwen3-4B. Grounding also increased mean input
tokens substantially: for example, Granite-4.1-3B rose from 128.1 to 627.8 tokens,
and Qwen3-4B from 128.4 to 760.0. These are per-assigned-query effects under the
frozen prompts, not intrinsic or token-normalized model-efficiency rankings.

\subsection{The Granite-4.0-1B Compatibility Failure}

Granite-4.0-1B scored zero both directly and grounded. Inspection of all clean A04,
A12, and A13 rows showed an identical 80-character sequence of exclamation marks;
every output reached the 80-token ceiling and none parsed. The A04 and A13 energy
endpoints also ran on different T4 boards. We therefore treat this pair as an
observed frozen prompt/chat-template/runtime incompatibility. It removed one
candidate from the useful online pool, but neither explains the policy-level
calibration infeasibility by itself nor shows that the model family is generally
incapable of RAG.

\subsection{Retrieval, Memory, and Failure Layers}

On the same board, A00 achieved 12.5\% while tiny BM25, dense, and hybrid routes
each achieved zero. Their measured energy increases were 102.97, 95.70, and
88.44 J/query, respectively, with intervals excluding zero. Again, the success
difference is the loss of no-retrieval cases rather than a decline from positive
evidence-task performance.

Tiny flat, hierarchical, and graph-memory routes also scored zero and showed
larger energy means than A00, but they ran on a different board; those energy
differences are descriptive. Across the common grounded route, retrieval recovered
the complete required evidence for most evidence-requiring queries, yet success
remained generator-dependent. For example, two-hop evidence was frequently
retrieved while Granite-4.1-3B and SmolLM3 still obtained zero strict success on
that class. Retrieval completeness was therefore necessary but not sufficient.

Failure decomposition reinforces this point. M18 parsed 96.67\% of outputs,
passed the support guard on 74.17\%, and succeeded on 44.17\%. M20 parsed 100\%,
passed support on 55.83\%, and succeeded on 34.17\%. A14 parsed 84.17\%, passed
support on 67.50\%, and succeeded on 42.50\%. In contrast, A03 parsed 46.67\% but
had zero support-qualified outputs. Sixty-four A03 outputs failed parsing and 40
contained citation representations that did not satisfy exact ID membership. The
frozen interface contract, not retrieval alone, was a major bottleneck.

\subsection{Robustness Is a Floor Effect}

\begin{table}[t]
\caption{Selected-policy outcomes under four evidence corruptions. Zero clean and corrupted success makes every success delta a floor effect.}
\label{tab:robustness}
\centering
\scriptsize
\begin{tabular}{@{}lrrr@{}}
\toprule
Condition & Success (\%) & Abstain (\%) & Coverage (\%) \\
\midrule
conflict injection & 0.0 & 15.8 & 100.0 \\
deletion duplicates & 0.0 & 15.0 & 100.0 \\
missing & 0.0 & 18.3 & 100.0 \\
stale & 0.0 & 15.0 & 100.0 \\
\bottomrule
\end{tabular}
\end{table}


The selected policy used A03 for every query under all four corruptions and achieved
zero strict success, matching its clean zero. The resulting zero degradation is
uninformative: it reflects a floor, not retained capability. A00 remained at 12.5\%
because it ignores evidence, while A13 remained at zero because of the compatibility
failure. The synthetic compromise marker was never emitted, but a zero marker rate
alongside zero policy efficacy cannot support a broad prompt-injection robustness
claim.
